% ===========================================================================
%  sviluppatore/04-motore/03-allocatore.tex
% ===========================================================================
\section{Dall'intenzione alla posizione}\label{sez:allocatore}

Una strategia dice \emph{che cosa}; l'allocatore decide \emph{quanto}. Sta in
\file{core/allocator.py}, e il suo comportamento si governa con una sola
\cod{dataclass} di configurazione.

\subsection{Le quattro modalità}

\begin{center}
\begin{tabular}{@{}L{28mm}L{110mm}@{}}
\toprule
\cod{strategy\_mix\_mode} & \textbf{Come distribuisce} \\
\midrule
\rigaTenue
\cod{equal\_weight}
  & Gerarchico classico: il portafoglio si divide fra le strategie, poi
    dentro ciascuna. Il fattore $\sqrt{N_\text{strat}}$ assume indipendenza
    fra strategie. È il valore predefinito, ed è anche il ripiego quando le
    altre modalità non hanno i dati che servono. \\
\cod{stream\_level}
  & Ogni strategia opera come stream indipendente con un proprio budget di
    volatilità (\cod{stream\_target\_vol\_annual}, \num{0.12} di
    default); i contributi si sommano con peso $1/N$. \\
\rigaTenue
\cod{explicit}
  & Si esegue il vettore di pesi dichiarato dalla strategia, punto. Serve a
    chi ha già calcolato le proprie dimensioni e non vuole che nessuno le
    tocchi. \\
\cod{covariance}
  & Le correlazioni si \emph{stimano} invece che assumerle: la volatilità del
    libro è $\sqrt{w^{\mathsf T}\Sigma w}$, con $\Sigma$ da
    \file{core/covariance.py}. \\
\rigaTenue
\cod{bayesian}
  & Dichiarata e non implementata: solleva \cod{NotImplementedError}. È
    elencata nei valori ammessi perché il messaggio d'errore è più utile di
    un «valore non valido». \\
\bottomrule
\end{tabular}
\end{center}

\begin{perche}{Perché \cod{explicit\_vol\_target} è un campo a sé}
In modalità \cod{explicit} si può chiedere che il libro venga scalato una
volta sola perché la sua volatilità prevista sia un certo valore --- le
proporzioni fra i titoli restano identiche, cambia solo quanto pesa il libro
nel suo insieme. Il parametro non è \cod{portfolio\_target\_vol\_annual}
perché quello ha sempre un valore (ha un default), e con un default non si
distingue «lo voglio» da «c'era». Qui l'assenza significa qualcosa:
\cod{None} vuol dire «esegui ciò che ti ho dato, senza scalare».
\end{perche}

\subsection{I tetti di rischio}

\begin{center}
\begin{tabular}{@{}L{44mm}R{16mm}L{76mm}@{}}
\toprule
\textbf{Parametro} & \textbf{Default} & \textbf{Che cosa limita} \\
\midrule
\rigaTenue
\cod{portfolio\_target\_vol\_annual} & \num{0.10} & Volatilità annua
  obiettivo del portafoglio. \\
\cod{max\_position\_weight} & \num{0.10} & $|w|$ del singolo strumento. \\
\rigaTenue
\cod{max\_asset\_class\_gross} & \num{0.20} & $\sum |w|$ dentro una classe
  di attività. \\
\cod{max\_gross\_leverage} & \num{1.5} & $\sum |w|$ totale. \\
\rigaTenue
\cod{min\_vol\_floor} & \num{0.05} & Pavimento sulla volatilità stimata. \\
\bottomrule
\end{tabular}
\end{center}

\begin{nota}{Il ragionamento dietro il tetto per classe}
Il commento nel codice lo dice senza giri: quattro obbligazioni tutte lunghe
fanno il \num{40}\,\% lordo senza alcun tetto, e in situazione di stress sono
una posizione sola. Il tetto al \num{20}\,\% limita a «un'esposizione
obbligazionaria con leva~2 al massimo», che è realistico. Non è un numero
scelto per prudenza generica: è la risposta a un modo preciso in cui una
diversificazione apparente sparisce.
\end{nota}

\subsection{La validazione è nel costruttore}

\simbolo{\_\_post\_init\_\_} controlla ogni campo e solleva \cod{ValueError}
con il valore ricevuto nel messaggio. Due vincoli meritano attenzione perché
non sono ovvi:

\begin{itemize}
  \item \cod{max\_gross\_leverage >= 1.0}. Un tetto inferiore a~\num{1}
        significherebbe non poter mai investire tutto il capitale, ed è quasi
        sempre un errore di unità di misura.
  \item \cod{max\_asset\_class\_gross} ammesso fino a~\num{4.0}. Il limite è
        stato ampliato per sostenere i portafogli lungo-corto, che arrivano
        a un lordo di~\num{2} e oltre.
\end{itemize}

\begin{attenzione}{La configurazione è congelata}
\simbolo{AllocatorConfig} è una \cod{dataclass} congelata: non si modifica un
campo dopo averla costruita. Per variarne uno in un test si usa
\cod{dataclasses.replace(config, max\_gross\_leverage=2.0)}, che ne
restituisce una copia nuova. Assegnare direttamente solleva
\cod{FrozenInstanceError}, ed è voluto: una configurazione di rischio
modificata a metà di un backtest renderebbe i risultati non riproducibili.
\end{attenzione}
